\chapter{Marco teorico}
\label{chap:marco}

\section{Opciones, futuros y volatilidad implicita}

Una opcion financiera otorga a su comprador el derecho, pero no la obligacion, de comprar o vender un activo subyacente en condiciones predefinidas. En una call, el derecho es comprar; en una put, vender. El strike $K$ fija el precio de ejercicio y el vencimiento determina el horizonte temporal del contrato. En el caso estudiado, las opciones se negocian sobre el IBEX y se relacionan con futuros mensuales que actuan como subyacente operativo.

El precio de una opcion depende de varias magnitudes: nivel del subyacente, strike, tiempo hasta vencimiento, tipo de interes, dividendos o carry, y volatilidad esperada. La volatilidad no entra como una variable observable directa. En la practica se invierte un modelo de valoracion: se busca la volatilidad que iguala el precio teorico al precio observado. Esa volatilidad se denomina volatilidad implicita y se interpreta como una representacion de mercado, no como una estimacion directa de volatilidad historica.

La volatilidad implicita tiene ventajas claras. Permite comparar opciones de precios absolutos muy distintos, estudiar la estructura por vencimientos y strikes, y detectar patrones como smiles y skews. Tambien tiene limitaciones: depende del modelo de valoracion usado, puede ser inestable en opciones iliquidas, y puede amplificar errores de precio, tiempo o tipo de interes. Por eso la construccion del dataset es parte central de la metodologia.

\section{Modelo Black-76}

Black-76 valora opciones europeas sobre futuros. Si $F$ es el precio del futuro, $K$ el strike, $r$ el tipo libre de riesgo, $T$ el tiempo hasta vencimiento y $\sigma$ la volatilidad, las formulas de call y put son:

\[
C=e^{-rT}\left(FN(d_1)-KN(d_2)\right),
\]

\[
P=e^{-rT}\left(KN(-d_2)-FN(-d_1)\right),
\]

donde:

\[
d_1=\frac{\ln(F/K)+\frac{1}{2}\sigma^2T}{\sigma\sqrt{T}}, \qquad
d_2=d_1-\sigma\sqrt{T}.
\]

El proyecto usa esta relacion en sentido inverso. Para cada operacion de opcion se conocen precio negociado, strike, tipo de opcion, vencimiento, tipo compuesto y futuro subyacente asociado. La volatilidad implicita es la raiz de la ecuacion:

\[
g(\sigma)=P_{Black76}(F,K,T,r,\sigma)-P_{mercado}=0.
\]

El problema numerico exige acotar el dominio de busqueda, tratar casos degenerados y descartar observaciones que no admiten una solucion coherente. La configuracion del repositorio fija un rango de volatilidad entre $10^{-8}$ y $10$, tolerancia $10^{-8}$ y filtros adicionales sobre tipo de operacion y lag maximo del subyacente. Estas decisiones no son detalles de implementacion menores: determinan que precios entran en el estudio.

\section{Superficie de volatilidad}

Si Black-76 asumiera volatilidad constante, todas las opciones sobre el mismo subyacente y vencimiento compartirian la misma volatilidad. El mercado real no se comporta asi. La volatilidad implicita suele depender de la moneyness y del vencimiento. La moneyness mide la posicion relativa del strike frente al subyacente; una forma comun y estable es la log-moneyness:

\[
m=\ln\left(\frac{K}{F}\right).
\]

En este trabajo tambien se considera la forward moneyness, que incorpora el ajuste temporal por tipo de interes. La superficie resultante permite observar smiles, alas y term structures. Modelar esa superficie con aprendizaje automatico equivale a aproximar una funcion no lineal con restricciones economicas implicitas: suavidad razonable, sensibilidad distinta por zonas y comportamiento interpretable ante cambios de vencimiento, strike o subyacente.

\section{Aprendizaje automatico supervisado}

El problema de modelado se plantea como regresion supervisada. Cada observacion contiene un vector de caracteristicas $X_i$ y una volatilidad implicita observada $y_i$. El modelo aprende una funcion $f$ y produce $\hat{y}_i=f(X_i)$. Las metricas principales son MAE, RMSE y $R^2$:

\[
MAE=\frac{1}{n}\sum_i |y_i-\hat{y}_i|,
\]

\[
RMSE=\sqrt{\frac{1}{n}\sum_i (y_i-\hat{y}_i)^2},
\]

\[
R^2=1-\frac{\sum_i (y_i-\hat{y}_i)^2}{\sum_i (y_i-\bar{y})^2}.
\]

El MAE mide el error medio en unidades de volatilidad y es robusto ante errores extremos. El RMSE penaliza mas los errores grandes y se usa como metrica principal de seleccion. El $R^2$ mide la proporcion de variabilidad explicada frente a un predictor constante.

\section{Familias de modelos}

La regresion lineal se usa como baseline. Su principal ventaja es la interpretabilidad: cada coeficiente resume una sensibilidad promedio condicionada a las demas variables. Su limitacion es evidente en una superficie de volatilidad no lineal: solo captura curvatura si se introducen transformaciones y terminos de interaccion.

Random Forest promedia arboles entrenados sobre muestras bootstrap y subconjuntos de variables. Es robusto, capta interacciones y no exige escalado de variables. Su prediccion es dificil de expresar en una formula compacta, pero puede analizarse con SHAP, arboles surrogate y diagnosticos de superficie.

XGBoost construye un ensamble secuencial de arboles. Cada arbol corrige errores residuales del conjunto anterior y la regularizacion controla complejidad. En datos tabulares suele ser competitivo, aunque su comportamiento tambien requiere herramientas de explicabilidad para entender dependencias y regiones de error.

Las redes neuronales secuenciales aproximan funciones no lineales mediante capas densas. Requieren escalado numerico y mecanismos de control como early stopping. La red inspirada en tensor-train introduce una estructura compacta para representar interacciones, con la intencion de reducir parametros y organizar la complejidad.

\section{Validacion temporal y sesgos}

En datos financieros, la validacion aleatoria puede producir conclusiones optimistas. Si observaciones futuras influyen en el entrenamiento, aparece lookahead bias. Si los hiperparametros se seleccionan mirando el conjunto de test, aparece data snooping. El proyecto evita estas fugas mediante separacion temporal: train precede a validation y validation precede a test. Ademas, se controla que un mismo contrato de opcion no aparezca simultaneamente en particiones distintas.

La exclusividad por contrato es relevante porque un contrato puede negociarse durante varios dias. Si un modelo ve un contrato en entrenamiento y despues se evalua sobre el mismo contrato en test, puede aprender caracteristicas idiosincraticas de ese contrato y sobrestimar su capacidad de generalizacion.

\section{Explicabilidad global y local}

La explicabilidad global busca responder que variables gobiernan el modelo en conjunto. SHAP permite descomponer una prediccion en una suma de contribuciones:

\[
f(x)=\phi_0+\sum_{j=1}^{p}\phi_j.
\]

El valor base $\phi_0$ representa la prediccion media de referencia y cada $\phi_j$ mide la contribucion atribuida a una variable. El repositorio usa un explicador por permutacion, adecuado para tratar como caja negra familias heterogeneas.

La explicabilidad local se centra en una prediccion concreta. Un waterfall SHAP muestra que variables empujan la prediccion hacia arriba o hacia abajo respecto al valor base. Los vecinos cercanos aportan otra dimension: indican si la muestra tiene soporte historico en el espacio de features o si se parece poco a las observaciones de entrenamiento.

\section{Respuesta funcional: ICE y ALE}

Las curvas ICE muestran como cambia la prediccion de muestras individuales cuando se modifica una variable y se mantienen las demas constantes. Son utiles para observar heterogeneidad: dos opciones pueden reaccionar de forma distinta a cambios de moneyness o vencimiento. Su riesgo es generar contrafactuales irreales cuando combina valores que no aparecen en mercado.

Las curvas ALE miden efectos acumulados locales y reducen parte del problema de extrapolacion al usar intervalos donde existen datos. Por ello son especialmente utiles cuando las variables estan correlacionadas. En este TFM, ICE y ALE se interpretan conjuntamente: ICE muestra dispersion individual; ALE resume tendencia media local.

\section{Modelos equivalentes}

Los arboles surrogate y la regresion simbolica aproximan las predicciones del modelo principal, no la volatilidad verdadera. Esta distincion es central. Un surrogate con baja fidelidad no debe usarse para explicar el modelo principal; un surrogate con alta fidelidad puede resumir reglas o formulas utiles, pero sigue heredando las limitaciones del estimador que imita.

El dashboard conserva metricas de fidelidad como RMSE, MAE y $R^2$ entre predicciones del surrogate y predicciones del modelo principal. Esta capa ayuda a comunicar modelos complejos ante usuarios tecnicos y no tecnicos sin presentar reglas simplificadas como certezas absolutas.

\section{Tipos de interes y descuento}

En Black-76 el tipo de interes entra como factor de descuento $e^{-rT}$. En periodos de tipos cercanos a cero o negativos, como buena parte del intervalo 2017--2022 en Europa, su impacto puede parecer secundario frente a moneyness y vencimiento. Aun asi, ignorarlo introduce una inconsistencia: dos opciones identicas en strike y futuro pero con vencimientos distintos no tienen el mismo descuento.

El proyecto combina EONIA y euro short-term rate. La transicion entre referencias se controla por fecha de corte y spread configurado. Esta parte es importante porque el tipo usado en el solver no es un simple valor puntual: debe ser compatible con el periodo hasta vencimiento. En la base final el tipo medio es negativo, lo que obliga a no asumir formulas simplificadas pensadas para entornos de tipos positivos.

\section{Moneyness y vencimiento como ejes financieros}

La moneyness resume donde se encuentra el strike respecto al subyacente. En opciones at-the-money, pequenos cambios en volatilidad suelen tener gran impacto en precio. En opciones muy dentro o fuera de dinero, la sensibilidad puede comportarse de manera distinta y depender mas de liquidez, vencimiento y precision del precio observado. Por ello, segmentar o diagnosticar por moneyness es esencial.

El vencimiento introduce otra dimension. A vencimientos cortos, el precio puede ser muy sensible a pequenos errores de tiempo o de subyacente. A vencimientos largos, la superficie puede incorporar expectativas mas amplias y menor densidad de operaciones. El producto de log-moneyness y raiz de vencimiento incluido en las features busca capturar que el efecto de alejarse de ATM no es igual en todos los horizontes.

\section{Errores y residuos en volatilidad}

Medir error en volatilidad implicita requiere cuidado. Un error absoluto de 0,05 puede ser moderado o elevado segun el nivel de volatilidad, vencimiento y uso operativo. No obstante, MAE y RMSE son adecuados como primera aproximacion porque trabajan en la misma escala que la variable objetivo. El RMSE penaliza errores grandes y, por tanto, es especialmente util para detectar modelos que fallan en regiones extremas.

El residual $y-\hat{y}$ conserva signo. Si es positivo, el modelo subestima la volatilidad implicita; si es negativo, la sobreestima. En un contexto de cobertura, ambos errores pueden tener implicaciones distintas. Una subestimacion sistematica en alas podria infravalorar riesgo de movimientos extremos; una sobreestimacion sistematica podria sesgar decisiones de precio o margen.

\section{Explicabilidad y gobernanza de modelos}

La explicabilidad en este TFM se aproxima a la gobernanza de modelos. No se limita a dibujar importancias. Un modelo gobernable debe permitir reconstruir datos, justificar seleccion de hiperparametros, comprobar rendimiento fuera de muestra, localizar fallos y comunicar limites. SHAP, ICE, ALE y surrogates son herramientas dentro de ese marco.

Tambien hay que reconocer sus riesgos. SHAP depende de una distribucion de referencia y de como se tratan dependencias entre variables. ICE puede mostrar contrafactuales sin soporte historico. ALE reduce ese problema, pero resume efectos medios locales. Un arbol surrogate puede parecer claro y aun asi tener fidelidad insuficiente. La regresion simbolica puede encontrar una formula compacta pero inestable fuera de muestra. Por eso el dashboard combina explicadores y conserva metricas de fidelidad.

\section{Complejidad frente a interpretabilidad}

La relacion entre complejidad e interpretabilidad no es lineal. Un modelo lineal es facil de leer, pero si no captura la superficie, sus coeficientes explican una aproximacion pobre. Un Random Forest es menos transparente internamente, pero si predice mejor y se acompana de buenos explicadores, puede ser mas util. La interpretabilidad relevante no es solo transparencia del algoritmo, sino capacidad de responder preguntas concretas con evidencia.

Este punto justifica comparar varias familias. La regresion lineal fija el minimo interpretable. Random Forest y XGBoost representan modelos tabulares potentes. Las redes prueban si una aproximacion diferenciable y parametrica aporta ventajas. La red tensor-train explora una estructura mas compacta para interacciones. La decision final se apoya en metricas, diagnosticos y explicabilidad.

\section{Riesgo de medicion en volatilidad implicita}

La variable objetivo del trabajo no procede de una etiqueta manual ni de una medicion fisica directa. Es el resultado de una inversion numerica. Por ello existe riesgo de medicion. Si el precio de mercado contiene ruido, si el futuro subyacente asociado no refleja exactamente el instante de la opcion, si el tipo compuesto no es representativo o si el tiempo hasta vencimiento se calcula con una convencion incorrecta, la volatilidad resultante puede desplazarse.

Este riesgo no invalida el estudio, pero condiciona su interpretacion. El modelo aprende la volatilidad implicita calculada bajo un conjunto de reglas. En consecuencia, una prediccion correcta significa coherencia con esa construccion, no necesariamente con una volatilidad ``verdadera'' independiente. La documentacion del ETL y de las validaciones es la forma de controlar este riesgo: cuanto mas explicitas son las reglas, mas auditable es la variable objetivo.

\section{Densidad de datos y extrapolacion}

Los modelos supervisados interpolan con mayor confianza en regiones densas del espacio de entrenamiento. En opciones, la densidad no es uniforme. Suele haber mas operaciones cerca de strikes y vencimientos liquidos, y menos en alas o vencimientos largos. Esto implica que la misma metrica global puede esconder zonas con soporte muy distinto.

La explicabilidad local mediante vecinos aborda esta cuestion. Si una muestra tiene vecinos cercanos en entrenamiento, la prediccion se apoya en informacion historica comparable. Si no los tiene, la prediccion puede ser extrapolacion. El dashboard no elimina esa extrapolacion, pero la hace visible. Esta visibilidad es especialmente importante en volatilidad, donde las zonas extremas pueden ser precisamente las mas relevantes para gestion de riesgo.
